\documentclass{article}
\usepackage[utf8]{inputenc}
\usepackage{hyperref}
\usepackage{listings}
\usepackage{xcolor}

\title{Academic Vibing: Structured Curiosity for Human-Agent Collaboration}
\author{Attogram \\ \url{https://github.com/attogram/academic-vibing}}
\date{October 2026}

\begin{document}

\maketitle

\begin{abstract}
[ROCK TALK]: Academic Vibing v0.4. Transition to Preprint. Automated CDI hardening. Grounding via 2025-26 Vibe Coding literature. Addressing the "Vibe Slop" crisis through adversarial friction.
\end{abstract}

\section{Introduction}
Academic Vibing v0.4 formalizes the methodology for a post-vibe-hangover research landscape. As identified in recent literature (Kumar 2026, Mims 2026), pure "vibe coding" without structural rigor leads to "Vibe Slop." This version introduces the Automated CDI Kit to quantify research signal.

\section{Methodology}
The methodology utilizes a Recursive Agent Consensus (RAC) network, now enhanced with automated Consensus Divergence Index (CDI) measurement.

\section{Implementation}
\begin{lstlisting}[language=Python, caption=CDI Calculation Logic]
def calculate_cdi(data):
    total_claims = len(data)
    disputed_claims = sum(1 for item in data if len(set(item['consensus'].values())) > 1)
    return disputed_claims / total_claims
\end{lstlisting}

\section{Future Work}
Zenodo DOI generation and arXiv submission.

\end{document}
